\subsection{Evaluación con usuarios}
\label{sec:oe5}

\textbf{Pregunta.} Las pruebas automáticas de la sección anterior verifican el código, pero no muestran si
un acudiente, sin que nadie le explique la aplicación, logra armar y pagar el kit de su hijo. Esta
evaluación lo responde. \textbf{Método.} Es la fase F5 del OE5, con un protocolo
preregistrado: el plan de análisis, con tareas, métricas y criterios de éxito, se fechó el 28 de octubre de
2026, antes de la primera sesión, para no ajustar los criterios
a los resultados. Cada sesión fue una prueba de usabilidad moderada con pensamiento en voz alta
\citep{lazar_research_2017}, sin instrucciones. Las tareas siguen la compra: crear la cuenta y elegir colegio
y grado (T1), declarar personaje y gasto (T2), esperar la generación (T3), revisar y elegir un kit (T4),
cambiar o quitar un producto (T5), pagar (T6) y encontrar el pedido (T7); T8, pedir un cambio al asistente,
es exploratoria y no se planteó a los participantes. T1, T4 y T6 son críticas.

\textbf{Métrica.} Cada tarea se codificó como éxito sin ayuda, con ayuda, fallo o falla del sistema, con
criterio estricto: cualquier indicación de la moderadora cuenta como ayuda. Además se definieron tres
criterios. C1: el SUS (System Usability Scale, cuestionario de diez preguntas con puntaje de 0 a 100) debe
superar 68, el valor medio de referencia. C2: la carga percibida, medida con el RTLX (Raw Task Load Index, de
0 a 100, donde más es más carga), debe ser menor con Tukit que en la última compra habitual, con un intervalo
que excluya el cero y un tamaño del efecto $d_z \ge 0{,}5$. C3: cada tarea crítica debe lograrse sin ayuda en
al menos 80\,\% de las sesiones. También se aplicó el NPS (Net Promoter Score, cuánto recomendarían la
aplicación), que se reporta sin criterio.

\textbf{Datos.} Participaron seis acudientes del Colegio Agustiniano Ciudad Salitre, el colegio piloto,
reclutados por conveniencia entre el 2 y el 3 de noviembre (cinco por videollamada, uno presencial). El pago
usó la pasarela en modo simulado. La primera sesión de ACU-04 se cortó al subir la lista, así que de ella
solo se codifica T1; su segunda grabación muestra el pago aprobado, y por eso se cuentan los seis pagos,
aunque T6 se codifica sobre cinco sesiones. Las sesiones se grabaron con consentimiento, se transcribieron
en el equipo del proyecto y un segundo codificador revisó cada hoja contra las mismas fuentes y dejó 78
correcciones que se incorporaron antes del análisis.

\textbf{Resultado.} La Tabla~\ref{tab:oe5-tareas} da el éxito sin ayuda por tarea con su intervalo de
Wilson. Los seis acudientes completaron el recorrido
de la cuenta a un pedido pagado; el pago (T6) se resolvió sin ayuda en cuatro de cinco sesiones, en una
mediana de 43 segundos, y el registro (T1) en cuatro de seis. El resultado que cambia el diseño es T4: solo
uno de cinco revisó los kits y eligió sin ayuda. Se repitieron tres problemas. Primero, dos acudientes no
encontraron el camino a «Ver mis kits» al terminar la espera. Segundo, no entendieron de dónde sale el kit
(«yo no sé cuánto vale en el Faro», el comercio con el que esa acudiente compara precios) ni qué significa «confiable». Tercero, las listas de 82 productos les
generaron «estrés visual». Hubo además dos fallas puntuales: el asistente flotante tapó el botón «Elegir» y el kit Fan mostró solo el personaje de la primera de dos hijas. Ningún acudiente
usó el asistente por su cuenta.

\begin{table}[htbp]
\centering
\small
\caption{Éxito sin ayuda por tarea en las sesiones con acudientes (n = 6; ACU-04 solo en T1).
Intervalo de Wilson del 95\,\% y mediana del tiempo de la tarea; en T3, la duración de la espera.}
\label{tab:oe5-tareas}
\begin{tabular}{@{}llccc@{}}
\toprule
Tarea & Qué pide & Sin ayuda & IC 95\,\% & Mediana (s) \\
\midrule
T1 (crítica) & Cuenta, colegio y grado & 4 de 6 (67\,\%) & 30--90\,\% & 70 \\
T2 & Personaje y opción de gasto & 3 de 5 (60\,\%) & 23--88\,\% & 32 \\
T3 & Esperar los kits & 4 de 5 (80\,\%) & 38--96\,\% & 50 \\
T4 (crítica) & Revisar los kits y elegir & 1 de 5 (20\,\%) & 4--62\,\% & 91 \\
T5 & Cambiar o quitar un producto & 3 de 5 (60\,\%) & 23--88\,\% & 84 \\
T6 (crítica) & Pagar (simulado) & 4 de 5 (80\,\%) & 38--96\,\% & 43 \\
T7 & Encontrar el pedido & 2 de 2 (100\,\%) & 34--100\,\% & 16 \\
\bottomrule
\end{tabular}
\end{table}

Los cuestionarios, que los seis respondieron el 5 de noviembre, resumen la valoración
(Tabla~\ref{tab:oe5-cuestionario}). El límite inferior del intervalo del SUS supera el umbral de 68 (prueba t
de una muestra, $p = 0{,}0004$) y en la escala de Bangor et al.~\citep{bangor_determining_2009} el promedio
corresponde a «excelente», un valor que orienta y no estima la población. La carga percibida con Tukit fue
mucho menor que la de la última compra habitual, con un intervalo de la diferencia que excluye el cero (prueba
t pareada, $p = 0{,}003$; $d_z$ es la diferencia media dividida por la desviación estándar de las
diferencias). Cuatro de seis describieron la aplicación como fácil, intuitiva o rápida («fue súper rápido»,
ACU-05), y valoraron el ahorro de tiempo y poder cambiar productos. El tiempo activo de las sesiones completas fue de 2,4 a 10,3 minutos, frente a cuatro horas o más
que reporta el 65\,\% de los acudientes en la encuesta del OE1; es solo un orden de magnitud, porque compara
una tarea guiada con una compra real.

\begin{table}[htbp]
\centering
\small
\caption{Cuestionario de cierre (n = 6). SUS y RTLX de 0 a 100; intervalos del 95\,\%
($t$ para SUS y la diferencia de RTLX, \emph{bootstrap} de 10.000 remuestreos para el NPS).}
\label{tab:oe5-cuestionario}
\begin{tabular}{@{}llcc@{}}
\toprule
Medida & Resultado & IC 95\,\% & Criterio \\
\midrule
SUS (C1) & 90 (de 77,5 a 100) & 82--98 & $> 68$: cumple \\
RTLX, compra habitual & 63 & & \\
RTLX, con Tukit & 9 & & \\
Diferencia RTLX (C2) & 54, $d_z = 1{,}83$ & 23--85 & excluye 0 y $d_z \ge 0{,}5$: cumple \\
NPS (S1) & $+83$ (5 promotores, 1 pasivo) & 50--100 & se reporta \\
\bottomrule
\end{tabular}
\end{table}

\textbf{Lectura.} El OE5 pide pruebas funcionales, no funcionales, de carga y de usabilidad con usuarios
reales. Las cuatro clases se ejecutaron, aunque las de carga y degradación quedaron en una corrida por equipo
y sin repetir todos los niveles (Sección~\ref{sec:calidad}). C1 y C2 se cumplen y C3 se alcanzó en el pago (T6),
pero no en el registro (T1) ni en la elección del kit (T4). Con seis participantes los resultados son
formativos \citep{nielsen_estimating_1994}, es decir, sirven para ver dónde se concentra la ayuda y orientar
la siguiente versión. Los primeros cambios apuntan a T4: llevar al acudiente a sus kits cuando estén listos, agrupar los productos por categoría con la tienda
visible, explicar las etiquetas de nivel y los textos de consolidación, corregir el personaje del segundo
hijo, reubicar el asistente y aclarar que la cuenta se crea al armar el kit.

\textbf{Límite.} El plan se desvió en varios puntos. El reclutamiento fue por conveniencia, con una sola ronda y una sola moderadora, y varios participantes conocen al equipo, así que el análisis se apoya en lo observado y no en los comentarios positivos. En el procedimiento, las
sesiones se planearon de 15 minutos, el cierre fue una entrevista oral y el SUS, el RTLX y el NPS se
respondieron por formulario dos o tres días después, de modo que dependen del recuerdo y de la cortesía; el
SEQ (una pregunta de dificultad después de cada tarea) no se aplicó. En la codificación, la moderadora
explicó más de lo que permitía el guion en dos sesiones, y esas ayudas cuentan como ayuda, lo que hace conservadoras las tasas. No se midió el conteo de interacciones (lo sustituyó el RTLX) ni la recodificación ciega del 30\,\% con kappa. Un acudiente pareció invertir dos escalas del SUS y su puntaje (77,5) se reporta tal cual. El pago simulado aprueba solo, así que T6 mide la navegación y no una pasarela real.
